%! AP Statistics — Unit 3, Lesson 3.2 — Guided Notes KEY
\documentclass[10pt]{article}
\usepackage{apstats-article}
\usepackage{apstats-key}

\begin{document}

\pageheader{Unit 3, Lesson 3.2}{Guided Notes — KEY}
\namedateperiod

\vspace{0.12in}

\begin{objectivebox}
\small By the end of this lesson, I will be able to\dots\\[4pt]
\ans{Distinguish between a population and a sample.}\\[1pt]
\ans{Classify a study as an observational study or an experiment.}\\[1pt]
\ans{Determine when generalizations and causal conclusions are valid based on the study type.}
\end{objectivebox}

\vspace{0.1in}

\begin{vocabbox}
\termblanklong{Population} \ans{The entire group of individuals about whom we want information.}
\termblanklong{Sample} \ans{A subset of the population selected for study.}
\termblanklong{Observational study} \ans{Researchers observe without imposing treatments.}
\termblanklong{Experiment} \ans{Researchers deliberately impose treatments on experimental units.}
\termblanklong{Generalizability} \ans{Whether conclusions from a sample can be extended to the population.}
\termblanklong{Causal relationship} \ans{One variable causes a change in another; requires an experiment with random assignment.}
\end{vocabbox}

\vspace{0.1in}

\begin{hookbox}
\small
\ans{The drug trial headline — researchers randomly assigned patients, so they can conclude the drug \emph{caused} the effect. The dog-owner headline is an observational study — an association, not causation (confounding variables could explain the link).}
\end{hookbox}

\vspace{0.1in}

\begin{notesbox}{Population vs. Sample}
\small
\renewcommand{\arraystretch}{1.6}
\begin{tabularx}{\linewidth}{@{} >{\bfseries\color{navy}}p{0.22\linewidth} X @{}}
Population & All \ans{individuals or units} of interest. \\
Sample & A \ans{subset} of the population selected for study. \\
\end{tabularx}

\vspace{10pt}
\textbf{Example:}
\begin{itemize}[itemsep=2pt]
  \item Population: \ans{All U.S. adults}
  \item Sample: \ans{The 500 randomly chosen adults surveyed}
\end{itemize}
\end{notesbox}

\vspace{0.1in}

\begin{notesbox}{Observational Study vs. Experiment}
\small
\renewcommand{\arraystretch}{1.6}
\begin{tabularx}{\linewidth}{@{} >{\bfseries\color{navy}}p{0.22\linewidth} X @{}}
Observational study & Researchers \ans{observe} without \ans{imposing} treatments. Includes surveys, retrospective and prospective studies. \\[4pt]
Experiment & Researchers \ans{impose} treatments to \ans{experimental} units. \\
\end{tabularx}

\vspace{10pt}
\textbf{Conclusion Rules (DAT-2.B):}
\begin{enumerate}[itemsep=4pt]
  \item Can only generalize to the \ans{population} from which the sample was \ans{randomly} selected.
  \item \textbf{Observational study} $\to$ \ans{cannot} establish causation (DAT-2.B.3).
  \item \textbf{Experiment with random assignment} $\to$ \ans{can} establish causation.
\end{enumerate}
\end{notesbox}

\vspace{0.1in}

\begin{practicebox}
\small
\vspace{6pt}
\renewcommand{\arraystretch}{1.8}
\begin{tabularx}{\linewidth}{@{} X p{2.4cm} p{2cm} p{2cm} @{}}
\textbf{Study} & \textbf{Type} & \textbf{Generalize?} & \textbf{Causation?} \\
\hline
Researchers randomly assign 100 patients to drug or placebo; measure recovery time. &
\ans{Experiment} & \ans{No} & \ans{Yes} \\
A survey of 500 randomly selected adults asks about screen time and sleep quality. &
\ans{Observ. study} & \ans{Yes} & \ans{No} \\
A teacher asks the 30 students in her class whether they like the new homework policy. &
\ans{Observ. study} & \ans{No} & \ans{No} \\
\end{tabularx}
\end{practicebox}

\end{document}
